% ============================================================
% CAPÍTULO 6
% ============================================================
\chapter{Despliegue e Instalación}

\section{Requisitos Previos}

Antes de iniciar el ensamblaje de FERMIER v7.1.6-R, verifique la disponibilidad de los siguientes elementos:

\begin{table}[h]
\centering
\begin{tabular}{|l|l|p{6cm}|}
\hline
\textbf{Categoría} & \textbf{Ítem} & \textbf{Especificación / Notas} \\ \hline
\multirow{5}{*}{Herramientas} 
 & Torquímetro dinamométrico & Rango 0.5--5 N·m, precisión $\pm$2\% \\ 
 & Llaves Allen métricas & 1.5, 2.0, 2.5, 3.0, 4.0 mm \\ 
 & Multímetro digital & True-RMS, capacidad de medición LiPo 12S (hasta 50 V) \\ 
 & Osciloscopio portátil (opcional) & Verificación de conmutación SSR ($<2\,\mu$s) \\ 
 & Estación de soldadura & Punta fina, temperatura regulada 320--350 $^\circ$C \\ \hline
\multirow{3}{*}{Software} 
 & Firmware de vuelo & ArduPilot / PX4 compatible con H-Frame asimétrico \\ 
 & Kernel PNLIO v2.1-OPT & Repositorio: \texttt{github.com/godear6959-creator/PNLIO-Framework} \\ 
 & Intérprete Python 3.10+ & Dependencias: \texttt{numpy}, \texttt{sqlite3}, \texttt{json}, \texttt{ollama} (modo local) \\ \hline
\multirow{3}{*}{Consumibles} 
 & Loctite 243 & Fijación de tornillería en aluminio 7075-T6 \\ 
 & Pasta térmica conductiva & Coeficiente $\geq$ 3 W/m·K para disipadores aviónica \\ 
 & Cinta de cobre EMI & Blindaje de la bahía de aviónica \\ \hline
\end{tabular}
\end{table}

\begin{quote}
\textbf{ADVERTENCIA:} La batería principal LiPo 12S 10000 mAh (44.4 V nominal) almacena 444 Wh de energía. Un cortocircuito puede generar corrientes superiores a 200 A. Utilice siempre guantes anticorte y gafas de protección durante la instalación eléctrica.
\end{quote}

\section{Ensamblaje de la Estructura Mecánica (H-Frame Asimétrico)}

El H-Frame asimétrico de FERMIER v7.1.6-R prioriza la estabilidad longitudinal para compensar el desplazamiento del centro de gravedad (CG) hacia la bahía de aviónica y la batería principal.

\subsection{Montaje del Chasis Central}

\begin{enumerate}
    \item \textbf{Placas centrales:} Fije las dos placas centrales de aluminio 7075-T6 (180 g cada una) utilizando los cuatro perfiles de unión de fibra de carbono. Aplique Loctite 243 en cada interfaz roscada.
    \item \textbf{Par de apriete:} 2.8 N·m en los tornillos M3$\times$12 de unión central; 2.2 N·m en los tornillos M2.5$\times$8 de los brazos.
    \item \textbf{Verificación de simetría:} Coloque el chasis sobre una superficie nivelada. La diferencia de altura entre los planos de motor M1--M2 y M3--M4 no debe exceder 1.0 mm. Corrija mediante arandelas de ajuste de 0.5 mm si es necesario.
\end{enumerate}

\subsection{Instalación de los Brazos y Motores}

Los brazos del H-Frame presentan una longitud diferencial: los brazos frontales (M1, M2) son 40 mm más cortos que los traseros (M3, M4) para contrarrestar la masa de la cámara NDVI montada en proa.

\begin{table}[h]
\centering
\begin{tabular}{|l|c|c|}
\hline
\textbf{Posición} & \textbf{Longitud del brazo} & \textbf{Sentido de rotación} \\ \hline
M1 (frontal izquierdo) & 220 mm & Horario (CW) \\ \hline
M2 (frontal derecho) & 220 mm & Antihorario (CCW) \\ \hline
M3 (trasero izquierdo) & 260 mm & Antihorario (CCW) \\ \hline
M4 (trasero derecho) & 260 mm & Horario (CW) \\ \hline
\end{tabular}
\end{table}

Asegure los motores brushless en las zapatas de aluminio. Verifique que el cableado de fases (A, B, C) permita una curvatura de radio $\geq$ 15 mm para evitar fatiga por vibración.

\subsection{Bahía de Aviónica y Blindaje EMI}

\begin{enumerate}
    \item Monte la bahía de aviónica en el compartimento central, desplazada 25 mm hacia la proa respecto al centro geométrico.
    \item Aplique cinta de cobre conductiva en el interior de la bahía, creando una jaula de Faraday continua con una apertura máxima de malla equivalente $< \lambda/20$ a 2.4 GHz.
    \item Instale la divisoria de Kevlar entre la aviónica y la batería principal. Esta divisoria actúa como barrera térmica (reducción de 12--15 $^\circ$C en la bahía durante vuelo sostenido) y como contención estructural ante eventos de inflado de la LiPo.
\end{enumerate}

\section{Instalación del Sistema Eléctrico}

\subsection{Batería Principal y Auxiliar}

\begin{table}[h]
\centering
\begin{tabular}{|l|l|l|}
\hline
\textbf{Parámetro} & \textbf{Batería Principal} & \textbf{Batería Auxiliar} \\ \hline
Ubicación & Compartimento inferior central & Compartimento superior trasero \\ \hline
Conexión & XT90-S antichispa & XT60 \\ \hline
Fusible & 80 A automático & 15 A rápido \\ \hline
Cableado principal & AWG 10, silicona & AWG 16, silicona \\ \hline
\end{tabular}
\end{table}

\begin{enumerate}
    \item Inserte la batería principal LiPo 12S 10000 mAh en su compartimento. El arnés de sujeción debe permitir extracción en $<$ 10 s para protocolos de emergencia.
    \item Conecte la batería auxiliar LiPo 12S 2200 mAh al módulo SSR. Este módulo monitorea la tensión de la principal en tiempo real.
\end{enumerate}

\subsection{Módulo SSR (Solid-State Switchover)}

El SSR constituye el núcleo de redundancia eléctrica de FERMIER. Su tiempo de conmutación garantizado es $<2\,\mu$s.

\textbf{Procedimiento de cableado:}
\begin{enumerate}
    \item \textbf{Entrada principal:} Conecte la salida del BEC de la batería principal (regulado a 5.2 V / 6 A) al pin \texttt{VIN\_MAIN} del SSR.
    \item \textbf{Entrada auxiliar:} Conecte la salida del BEC de la batería auxiliar al pin \texttt{VIN\_AUX}.
    \item \textbf{Salida aviónica:} El pin \texttt{VOUT} alimenta el flight controller, la cámara NDVI y el microcontrolador del kernel PNLIO.
    \item \textbf{Umbral de disparo:} El firmware del SSR está calibrado de fábrica a \textbf{39.6 V} (corregido en v7.1.6-R, vs.\ 15.0 V en v7.1.5-R). Este umbral protege contra descarga profunda de las celdas LiPo 12S ($3.3\,\text{V/celda} \times 12 = 39.6\,\text{V}$).
\end{enumerate}

\textbf{Verificación del SSR:}
\begin{itemize}
    \item Aplique tensión nominal (44.4 V) en \texttt{VIN\_MAIN}. Mida \texttt{VOUT} = 5.2 V.
    \item Reduzca \texttt{VIN\_MAIN} lentamente con una carga resistiva. Al alcanzar 39.5 V, el LED de estado del SSR debe cambiar de verde a ámbar; a 39.6 V, la conmutación a \texttt{VIN\_AUX} debe ocurrir sin caída de tensión detectable en \texttt{VOUT} (verificar con osciloscopio, caída $<$ 50 mV).
\end{itemize}

\subsection{Panel Solar y MPPT 12S}

El módulo solar de 25 W (monocristalino, eficiencia $\geq$ 20\%) se instala en la cubierta superior del H-Frame.

\begin{enumerate}
    \item Fije el panel con cuatro soportes articulados de nylon reforzado. El ángulo de incidencia óptimo para carga en terreno es 30$^\circ$ respecto a la horizontal.
    \item Conecte la salida del panel (18 V / 1.4 A en condiciones STC) al regulador MPPT 12S.
    \item El MPPT eleva la tensión a 50.4 V ($4.2\,\text{V/celda} \times 12$) para carga de mantenimiento de la batería auxiliar. El tiempo de carga completa desde 20\% SOC es de 6--7 horas de insolación directa.
\end{enumerate}

\begin{quote}
\textbf{Nota de versión v7.1.6-R:} El panel solar y el MPPT 12S son adiciones respecto a v7.1.5-R. Anteriormente, la batería auxiliar solo se recargaba en base. Esta mejora permite autonomía extendida en operaciones de terreno.
\end{quote}

\section{Instalación de la Aviónica}

\subsection{Flight Controller y Sensores}

\begin{enumerate}
    \item Monte el flight controller sobre amortiguadores de gel de sílice (duro 30 Shore A) para aislamiento de vibraciones de los motores ($>$ 80 Hz).
    \item Oriente la flecha de proa del FC coincidente con el eje longitudinal del H-Frame (dirección M1--M2).
    \item Conecte los ESCs a los puertos de salida siguiendo la secuencia: M1 $\rightarrow$ Salida 1, M2 $\rightarrow$ Salida 2, M3 $\rightarrow$ Salida 3, M4 $\rightarrow$ Salida 4.
\end{enumerate}

\subsection{Cámara NDVI}

La cámara multiespectral NDVI se monta en la proa, suspendida en una plataforma de amortiguación de caucho butílico.

\begin{itemize}
    \item \textbf{Interfaz:} UART (telemetría) + GPIO (trigger).
    \item \textbf{Resolución:} 1280$\times$960 px, banda roja (660 nm) e infrarroja cercana (840 nm).
    \item \textbf{Calibración radiométrica:} Realice una calibración con panel reflectante blanco (Spectralon) antes del primer vuelo. El coeficiente de calibración se almacena en el archivo \texttt{cal\_ndvi.json} del kernel PNLIO.
\end{itemize}

\subsection{Unidad de Procesamiento del Kernel PNLIO}

El kernel PNLIO v2.1-OPT requiere una unidad de procesamiento local. La configuración mínima validada es:

\begin{table}[h]
\centering
\begin{tabular}{|l|p{8cm}|}
\hline
\textbf{Componente} & \textbf{Especificación} \\ \hline
Microcontrolador / SBC & Raspberry Pi 4B (4 GB RAM) o equivalente ARM64 \\ \hline
Almacenamiento & MicroSD UHS-I 64 GB (clase A2) para SQLite WAL \\ \hline
Comunicación & UART al flight controller, I2C a sensores ambientales \\ \hline
Runtime & Ollama 0.3.x + modelo Phi-3 Mini (modo air-gapped) \\ \hline
\end{tabular}
\end{table}

Instale el kernel en \texttt{/opt/pnlio/}. La estructura de directorios debe ser:

\begin{verbatim}
/opt/pnlio/
├── kernel/
│   ├── pnlio_core.py
│   ├── metrics.py          # Δθ, C, IRE
│   └── ssr_guard.py        # Monitoreo de conmutación
├── data/
│   ├── agronomic_vector.json
│   ├── calibration/
│   └── telemetry_wal.db
└── models/
    └── phi3_q4.gguf
\end{verbatim}

\section{Carga de Firmware y Calibración Inicial}

\subsection{Firmware de Vuelo}

\begin{enumerate}
    \item Cargue el firmware ArduPilot (Copter 4.5+) o PX4 (v1.14+) compatible con H-Frame.
    \item Configure el tipo de frame: \texttt{FRAME\_CLASS = 3} (H-Frame) en ArduPilot, o \texttt{CA\_AIRFRAME = 1} (quad H) en PX4.
    \item Ajuste los parámetros de asimetría: \texttt{MOT\_YAW\_HEADROOM} incrementado en 15\% para compensar el brazo trasero más largo.
\end{enumerate}

\subsection{Calibración del Kernel PNLIO}

Ejecute el script de inicialización:

\begin{verbatim}
cd /opt/pnlio/kernel
python3 init_pnlio.py --calibrate
\end{verbatim}

Este script realiza:
\begin{enumerate}
    \item \textbf{Calibración de $\Delta\tau$:} Sincroniza el reloj de alta resolución del SBC con \texttt{time.time()} (corrección v7.1.6-R; en v7.1.5-R se usaba un contador de ciclos propenso a desbordamiento).
    \item \textbf{Carga del vector agronómico:} Lee \texttt{agronomic\_vector.json} y valida su esquema contra el esquema PNLIO v2.1. En versiones anteriores, este vector se generaba aleatoriamente, introduciendo incertidumbre en las métricas.
    \item \textbf{Prueba de SQLite WAL:} Crea la base de datos \texttt{telemetry\_wal.db} con modo WAL habilitado, garantizando integridad de datos en escrituras concurrentes durante vuelo.
\end{enumerate}

\subsection{Checklist Pre-Instalación}

\begin{table}[h]
\centering
\begin{tabular}{|c|p{10cm}|c|}
\hline
\textbf{\#} & \textbf{Verificación} & \textbf{Estado} \\ \hline
1 & Tornillería central apretada a torque especificado & $\square$ \\ \hline
2 & CG dentro del 5\% del centro geométrico & $\square$ \\ \hline
3 & SSR conmuta correctamente a 39.6 V & $\square$ \\ \hline
4 & Tensión de batería principal $\geq$ 44.0 V ($\geq$ 3.67 V/celda) & $\square$ \\ \hline
5 & Tensión de batería auxiliar $\geq$ 44.0 V & $\square$ \\ \hline
6 & Cámara NDVI responde a trigger GPIO & $\square$ \\ \hline
7 & Kernel PNLIO carga vector JSON sin errores & $\square$ \\ \hline
8 & Modo air-gapped: sin interfaces WiFi/Bluetooth activas & $\square$ \\ \hline
9 & Panel solar limpio y sin microgrietas & $\square$ \\ \hline
10 & SQLite WAL activo y espacio libre en SD $>$ 10 GB & $\square$ \\ \hline
\end{tabular}
\end{table}


% ============================================================
% CAPÍTULO 7
% ============================================================
\chapter{Operación y Flujos de Trabajo}

\section{Preparación de Misión}

\subsection{Planificación del Vuelo}

FERMIER v7.1.6-R opera en entornos \textbf{air-gapped}; por tanto, toda la planificación se realiza previamente en estación base y se carga vía cable USB-C o tarjeta SD.

\begin{enumerate}
    \item \textbf{Definición del polígono:} Utilice QGroundControl o Mission Planner para delimitar el área agronómica. Exporte el archivo \texttt{.waypoints} al directorio \texttt{/opt/pnlio/missions/}.
    \item \textbf{Carga del vector agronómico:} El archivo \texttt{agronomic\_vector.json} debe contener, como mínimo, los siguientes campos estructurados:
\end{enumerate}

\begin{verbatim}
{
  "crop_type": "trigo_durum",
  "growth_stage": "espigazon",
  "expected_ndvi_range": [0.45, 0.82],
  "alert_thresholds": {
    "ndvi_min": 0.35,
    "ire_min": 2.1
  },
  "temporal_resolution_sec": 5.0
}
\end{verbatim}

El kernel PNLIO utiliza este vector para contextualizar las métricas matemáticas ($\Delta\theta$, $C$, $IRE$) en términos agronómicos concretos.

\subsection{Secuencia de Encendido}

Siga estrictamente el orden de encendido para evitar picos de corriente en la aviónica:

\begin{table}[h]
\centering
\begin{tabular}{|c|l|l|}
\hline
\textbf{Paso} & \textbf{Acción} & \textbf{Indicador Esperado} \\ \hline
1 & Insertar batería auxiliar & LED SSR: verde (modo standby) \\ \hline
2 & Insertar batería principal & LED SSR: verde brillante; FC inicia boot \\ \hline
3 & Esperar 15 s & GPS: fix 3D (mínimo 8 satélites) \\ \hline
4 & Activar switch de seguridad & Buzzer: tono de confirmación \\ \hline
5 & Iniciar kernel PNLIO & Terminal: \texttt{PNLIO v2.1-OPT ready. Δτ=time.time()} \\ \hline
6 & Verificar modo air-gapped & \texttt{ip link show} $\rightarrow$ sin interfaces wlan0, bluetooth0 \\ \hline
\end{tabular}
\end{table}

\section{Flujo de Validación en Tiempo Real (Kernel PNLIO)}

Durante el vuelo, el kernel PNLIO ejecuta un ciclo de validación matemática continua sobre cada paquete de telemetría.

\subsection{Adquisición de Datos Brutos}

El flight controller transmite a 10 Hz vía UART:

\begin{itemize}
    \item Vector de actitud: $\vec{v}_a$ (roll, pitch, yaw en radianes)
    \item Vector de heading: $\vec{v}_h$ (dirección de deriva del viento + rumbo GPS)
    \item Timestamp de alta resolución: $t$ (segundos desde epoch, float64)
\end{itemize}

\subsection{Cálculo de Métricas PNLIO}

Para cada muestra $i$, el kernel computa:

\textbf{1. Coherencia Angular ($\Delta\theta$):}
\begin{equation}
\Delta\theta_i = \vec{v}_{h,i} \cdot \vec{v}_{a,i}
\end{equation}
Donde $\cdot$ denota el producto punto normalizado. $\Delta\theta \in [-1, 1]$. Valores cercanos a 1 indican alineación perfecta entre heading y actitud.

\textbf{2. Coherencia Pura ($C$):}
\begin{equation}
C_i = \frac{\Delta\theta_i}{\max(10^{-6}, \Delta\tau_i)}
\end{equation}
Donde $\Delta\tau_i = t_i - t_{i-1}$ (diferencia de tiempo real, no contador de ciclos). El denominador protegido con $10^{-6}$ evita división por cero en congelamientos del sistema operativo.

\textbf{3. Entropía Textual ($H_{txt}$):}
$H_{txt}$ se deriva del análisis de la imagen NDVI recién capturada. Se cuantiza la distribución de píxeles en 256 bins y se aplica entropía de Shannon:
\begin{equation}
H_{txt} = -\sum_{k=0}^{255} p_k \log_2(p_k)
\end{equation}
Donde $p_k$ es la frecuencia normalizada del bin $k$.

\textbf{4. Índice de Realidad Extendida ($IRE$):}
\begin{equation}
IRE_i = C_i \cdot \log_2(1.0 + H_{txt})
\end{equation}
$IRE$ combina la estabilidad cinemática del dron ($C$) con la complejidad informativa de la escena agronómica ($H_{txt}$). Un $IRE$ alto indica datos fiables y ricos en información.

\subsection{Reglas de Validación y Almacenamiento}

\begin{table}[h]
\centering
\begin{tabular}{|p{4cm}|p{9cm}|}
\hline
\textbf{Condición} & \textbf{Acción del Kernel} \\ \hline
$C \geq 0.75$ (Efecto Reflex) & Datos validados. Se habilita almacenamiento en SQLite WAL y trigger de cámara NDVI para siguiente waypoint. \\ \hline
$0.50 \leq C < 0.75$ & Datos almacenados con flag \texttt{unstable}. Se emite advertencia en telemetría local. \\ \hline
$C < 0.50$ & Datos descartados. Se registra evento en \texttt{pnlio\_anomalies.log}. Se activa modo de espera de 2 s para nueva muestra. \\ \hline
$IRE <$ \texttt{alert\_thresholds.ire\_min} & Se marca el waypoint para reinspección. El flight controller reduce velocidad de crucero al 60\%. \\ \hline
\end{tabular}
\end{table}

\subsection{Persistencia Air-Gapped}

Toda la telemetría validada se escribe en \texttt{telemetry\_wal.db} utilizando el modo \textbf{Write-Ahead Logging (WAL)} de SQLite. Esta configuración ofrece:

\begin{itemize}
    \item \textbf{Atomicidad:} Las transacciones de telemetría son atómicas incluso ante corte de energía (mitigado por el SSR).
    \item \textbf{Concurrencia:} Lecturas de análisis (post-vuelo) no bloquean escrituras en vuelo.
    \item \textbf{Portabilidad:} La base de datos es un único archivo de $\sim$50 MB/h de vuelo, extraíble físicamente.
\end{itemize}

\section{Modos de Operación}

\subsection{Modo Vuelo Autónomo (MVA)}

\begin{itemize}
    \item \textbf{Descripción:} FERMIER ejecuta la misión preprogramada de waypoints.
    \item \textbf{Velocidad de crucero:} 8--12 m/s (ajustable por $IRE$ en tiempo real).
    \item \textbf{Altitud de trabajo:} 30--80 m AGL (Above Ground Level), óptima para NDVI.
    \item \textbf{Autonomía:} 30--40 minutos con batería principal al 100\% SOC.
\end{itemize}

\textbf{Flujo de trabajo MVA:}
\begin{verbatim}
INICIO → Carga waypoints → Despegue RTL → Navegación waypoint N
   ↓                                    ↓
Validación PNLIO ←── IRE ≥ umbral ──── Captura NDVI
   ↓                                    ↓
Almacena WAL ←── C ≥ 0.75 ─────────── Continúa a N+1
   ↓
FIN MISIÓN → RTL (Return to Launch)
\end{verbatim}

\subsection{Modo Terreno Solar (MTS)}

Activado manualmente tras aterrizaje o por fail-safe de batería principal $< 39.6$ V.

\begin{itemize}
    \item \textbf{Descripción:} FERMIER permanece en tierra como estación de monitoreo fija.
    \item \textbf{Alimentación:} Batería auxiliar + panel solar 25 W vía MPPT 12S.
    \item \textbf{Funciones:} Toma periódica de NDVI (cada 30 min), registro de datos ambientales (temperatura, humedad, radiación), transmisión local opcional vía LoRa 868 MHz.
    \item \textbf{Autonomía:} 6--7 horas continuas con insolación directa; 4 horas en condiciones nubladas parciales.
\end{itemize}

\subsection{Modo Air-Gapped de Emergencia (MAE)}

Activado automáticamente si el kernel detecta intento de handshake en puertos de red o si el SSR registra más de 3 conmutaciones en 60 s (indicativo de inestabilidad eléctrica severa).

\begin{itemize}
    \item Se apagan todos los transceptores RF no esenciales.
    \item El flight controller opera en modo estabilizado manual (si hay operador) o RTL autónomo.
    \item El kernel PNLIO entra en modo \textbf{logging mínimo}: solo almacena $\Delta\theta$, $C$ y estado del SSR, descartando $H_{txt}$ para reducir escrituras en SD.
\end{itemize}

\section{Protocolos de Emergencia}

\begin{table}[h]
\centering
\begin{tabular}{|p{3.5cm}|p{5cm}|p{5cm}|}
\hline
\textbf{Escenario} & \textbf{Condición de Disparo} & \textbf{Acción Automática} \\ \hline
Fail-safe de batería & $V_{\text{main}} < 39.6$ V durante $> 500$ ms & SSR conmuta a auxiliar en $<2\,\mu$s; FC ejecuta RTL inmediato \\ \hline
Pérdida de coherencia & $C < 0.25$ durante 5 s consecutivos & Hover en posición actual; espera recuperación de GPS/IMU \\ \hline
Pérdida de señal RC & No heartbeat en 3 s & RTL autónomo \\ \hline
Geofence violado & Posición fuera de polígono $\pm$ 20 m & Aterrizaje forzoso en punto seguro más cercano \\ \hline
Temperatura bahía & $T > 65^\circ$C & Reducción de frecuencia de muestreo PNLIO al 50\%; ventilación pasiva \\ \hline
\end{tabular}
\end{table}

\section{Post-Vuelo: Extracción y Análisis}

\subsection{Extracción de Datos}

\begin{enumerate}
    \item Aterrice FERMIER en el punto de despegue.
    \item Desconecte la batería principal (primero) y luego la auxiliar.
    \item Extraiga la microSD del SBC o conecte cable USB-C.
    \item Copie los archivos:
    \begin{itemize}
        \item \texttt{telemetry\_wal.db} (telemetría validada)
        \item \texttt{pnlio\_anomalies.log} (eventos de descarte)
        \item \texttt{ndvi\_captures/} (imágenes con metadatos $IRE$ incrustados en EXIF)
    \end{itemize}
\end{enumerate}

\subsection{Análisis RCR v2.0}

El índice de \textbf{Reducción de Confusión y Ruido (RCR) v2.0} cuantifica la calidad global de la misión:
\begin{equation}
RCR = \frac{N_{\text{valid}}}{N_{\text{total}}} \times \frac{\overline{IRE}}{\max(IRE)} \times \left(1 - \frac{N_{\text{anomalies}}}{100}\right)
\end{equation}
Donde:
\begin{itemize}
    \item $N_{\text{valid}}$ = muestras con $C \geq 0.75$
    \item $N_{\text{total}}$ = muestras totales
    \item $N_{\text{anomalies}}$ = eventos registrados en \texttt{pnlio\_anomalies.log}
\end{itemize}

Un $RCR > 0.85$ indica una misión de alta fidelidad, con reducción de alucinaciones/ruido del 87--92\% respecto a sistemas sin validación PNLIO.

\subsection{Mantenimiento Post-Misión}

\begin{table}[h]
\centering
\begin{tabular}{|l|c|p{7cm}|}
\hline
\textbf{Componente} & \textbf{Intervalo} & \textbf{Acción} \\ \hline
Tornillería estructural & Cada 5 vuelos & Retorqueo a especificación \\ \hline
Bahía de aviónica & Cada 10 vuelos & Inspección visual de integridad del blindaje de cobre \\ \hline
Conectores XT90/XT60 & Cada 3 vuelos & Limpieza con alcohol isopropílico; verificación de retención \\ \hline
Panel solar & Cada vuelo & Limpieza con paño microfibra; verificación de voltaje en circuito abierto ($V_{oc} \approx 21$ V) \\ \hline
Baterías LiPo & Cada vuelo & Almacenamiento a 3.85 V/celda (50\% SOC) si no se usarán en 48 h \\ \hline
MicroSD & Cada 20 vuelos & Formateo completo (exFAT) para prevenir fragmentación WAL \\ \hline
\end{tabular}
\end{table}

% ============================================================
% REFERENCIAS CRUZADAS (Opcional, para tu control)
% ============================================================
% Umbral SSR 39.6 V        -> Cap. 4, Tabla 4.1, ítem 3
% Δτ = time.time()         -> Cap. 4, Tabla 4.1, ítem 1
% Vector agronómico JSON   -> Cap. 4, Tabla 4.1, ítem 2
% Panel solar 25W + MPPT   -> Cap. 4, Tabla 4.1, ítem 4
% Estructura H-Frame       -> Cap. 2.1, Fig. 2.1, Tabla 5.1
% Ecuaciones Δθ, C, IRE    -> Cap. 3.1
% Efecto Reflex (C≥0.75)   -> Cap. 8, Glosario
